% !PRINTED 115
\weberpage{115}
Ferner ergibt sich aus (5) und den beiden letzten Gleichungen (9):

\weq{14}{f_1\!\left(-\frac{1}{\omega}\right) = f_2(\omega), \qquad f_2\!\left(-\frac{1}{\omega}\right) = f_1(\omega),}
und wenn man hiervon in der zweiten Gleichung (11) Gebrauch macht:

\weq{15}{f\!\left(-\frac{1}{\omega}\right) = f(\omega).}
Für die \gesperrt{Transformation zweiter Ordnung} folgt unmittelbar aus den beiden letzten Gleichungen (9):

\weq{16}{f_1(2\,\omega) f_2(\omega) = \sqrt{2},}
woraus sich noch durch Anwendung von (12) ergibt:

\weq{17}{f_2(\omega)^4\left[f(2\,\omega)^8 + f_2(2\,\omega)^8\right] = 2\left[f(\omega)^8 + f_1(\omega)^3\right].}
% [sic] last exponent printed as 3, inconsistent with the surrounding exponents of 8
Setzt man in (16) nach (13), (14)
\[
\begin{aligned}
f_1(2\,\omega) &= e^{-\frac{\pi i}{24}} f(2\,\omega - 1),\\
f_2(\omega) &= e^{\frac{\pi i}{24}} f\!\left(1 - \frac{1}{\omega}\right),
\end{aligned}
\]
so ergibt sich:
\[
f\!\left(1 - \frac{1}{\omega}\right) f(2\,\omega - 1) = \sqrt{2},
\]
und indem man $2\,\omega - 1$ gleich einem neuen $\omega$ setzt:

\weq{18}{f(\omega) f\!\left(\frac{\omega - 1}{\omega + 1}\right) = \sqrt{2}.}
Ersetzt man in (16) $\omega$ durch $\dfrac{\omega}{2}$ und $\dfrac{\omega+1}{2}$, so folgt:

\weq{19}{\begin{aligned}
f_1(\omega) f_2\!\left(\frac{\omega}{2}\right) &= \sqrt{2}, \\
f(\omega) f_2\!\left(\frac{\omega + 1}{2}\right) &= e^{\frac{\pi i}{24}} \sqrt{2}.
\end{aligned}}
Wir stellen endlich noch für ein ungerades $n$ die Formeln für die Haupttransformation $n$ter Ordnung der Funktionen $\eta$, $f$ auf. Für $\eta(n\,\omega)$ ergibt sich leicht, wenn man in den Formeln (6), \S\,33 nach der Differentiation $u = 0$ setzt, und die dritte Wurzel zieht:

\weq{20}{\sqrt{n}\,\eta(n\,\omega)\,\eta(\omega)^{\frac{n-3}{2}} = \prod_{1,\,\frac{n-1}{2}}^{\nu} \vartheta_{11}\!\left(\frac{\nu}{n}\right).}
